\scanpage{91}
% archive foliation: 90
\uline{Proposition} Soit $f:X-Y$ \deleted{morphisme} application constructible d'ens schématiques. Pour tte partie constructible $E$ de $Y$, $f^{-1}(E)$ est constructible.

\margin{\diagram{\begin{tikzcd}
\deleted{X}\arrow[r]\deleted{}&\deleted{Y}\\
\deleted{X_i}\arrow[u]\arrow[ur]&
\end{tikzcd}}}

En effet, soit $X=\coprod_iX_i$, les $X_i$ loc. fermés et constructibles dans $X$, \deleted{Alors} $f_i=(f\,|\,|X_i|)$ un morphisme. On a $f^{-1}(E)=\bigcup f_i^{-1}(\unsure{E_f|E})$, il suffit de prouver que chaque $f^{-1}(E)=f^{-1}(E)\cap X_i$ est constructible dans $X$, ou, ce qui revient au même, dans $X_i$ [\deleted{par} par le th. de Chevalley, \inserted{l'immersion} $X_i-X$ étant de présentation finie, pour une structure \unsure{induite} convenable sur $X_i$ \ldots]. Car on se ramène au cas d'un \uline{morphisme} \inserted{de préschémas}, comme il est quasi-compact quasi-séparé c'est connu.

\margin{\unsure{topologique} \unsure{l'argument}}

\uline{Corollaire} Invariance de la notion de partie constructible d'un ens schématique.

\uline{Proposition} Soit $f:X\longrightarrow Y$ une application constructible de \deleted{ens. schématiques} \inserted{préschémas quasi-compacts quasi-séparés}. Pour que $f$ soit un isom. d'ens schém., il faut et s. qu'il existe \deleted{\unsure{un}} pour un fini $I$, deux partitions $(X_i)_{i\in I}$ et $(Y_j)_{j\in I}$ de $X$ ($Y$) en parties loc. fermées constructibles \inserted{[\unsure{Affines si on veut}]}, telles que $f$ induise un \uline{isom} \uline{de} $X_i$ \uline{sur} $Y_i$.

C'est évidemment suffisant, prouvons la nécessité. Soit $g$ l'inverse de $f$, $X_\alpha,Y_\beta$ deux partitions \inserted{de $X,Y$} mettant en évidence le caractère constructible de $f,g$. \unsure{Évidem.} $f,g$ sont des bijections \unsure{inv.} l'une de l'autre.

\margin{\diagram{\begin{tikzcd}
X\arrow[r,shift left,"f"]&Y\arrow[l,shift left,"g"]
\end{tikzcd}}}
% The reversed blue typewriting is transcribed on its printed face, scan 92.
